\documentclass[./main.tex]{subfiles}


\begin{document}
\onehalfspacing
\section{Տոպոլոգիական տարածությունների անընդհատ արտա\-պատ\-կե\-րում\-ներ, թեորեմ մետրիկային տարածությունների արտա\-պատ\-կե\-րում\-ների անընդհատության մասին, անընդհատության հայտա\-նիշ\-ներ։ Հոմեոմորֆիզմ և հոմեոմորֆ տարածություններ, տոպո\-լո\-գիա\-կան տիպ և տոպոլոգիական հատկություններ։ Բաց, փակ արտա\-պատ\-կե\-րում\-ներ։}
\label{sec:9}

\par Մինչև այժմ մենք դիտարկում էինք տոպոլոգիական տարածություններն առան\-ձին-առանձին, միմյանցից անկախ։ Այժմ զբաղվելու ենք դրանց համեմատմամբ։ Այդ նպատակով ներմուծվում է տոպոլոգիական տարածությունների անընդհատ արտապակերման հասկացությունը, որը երկրորդ կարևորագույն հասկացությունն է տոպոլոգիական տարածություն հասկացությունից հետո։

\begin{definition}
Դիցուք ունենք $(X, \tau)$, $(Y, \sigma)$ տոպոլոգիական տարածություններ և ${f:X \rightarrow Y}$ արտապատկերում։ Ապա $f$-ը կոչվում է \textbf{անընդհատ $\boldsymbol{x_0\in X}$ կետում}, եթե $Y$ տարածությունում $f(x_0)=y_0$ կետի ամեն մի $V$ շրջակայքի համար գոյություն ունի $X$ տարածությունում $x_0$ կետի $U$ շրջակայք, որ $f(U)\subset V$։

\end{definition}
\par Հետևյալ պնդումը երբեմն հեշտացնում է անընդհատության պայմանի ստուգումը։

\begin{theorem}\label{ch9thm1} 
Ունենք տոպոլոգիական տարածությունների ${f:X \rightarrow Y}$ արտապատ\-կերում, $f(x_0)=y_{0}$, իսկ $\beta_{x_0}=\{U_i(x_0);\ i\in I\}$ և $\beta_{y_0}=\{V_j(y_0);\ j\in J\}$ ընտանիքները համապատասխանաբար $x_0$ և $y_0$ կետերի շրջակայքերի որևէ բազաներ են $X$ և $Y$ տարածություններում։ Ապա $f$-ը անընդհատ է $x_0$ կետում այն և միայն այն դեպքում, երբ $\forall\, V_j(y_0)\in \beta_{y_0}$ շրջակայքի համար գոյություն ունի $U_i(x_0)\in \beta_{x_0}$ շրջակայք, որ $f(U_i (x_0 ))\subset V_j(y_0)$։
\label{թեորեմ $1$}
\end{theorem}

\begin{proof} Դիցուք $f$-ը անընդհատ է $x_0$ կետում, և ունենք որևէ $V_j(y_0)\in \beta_{y_0}$։ Ըստ կետում անընդհատության սահմանման, գոյություն ունի $x_0$ կետի $U$ շրջակայք, որ $f(U)\subset V_j(y_0)$։ Համաձայն կետի շրջակայքերի բազայի սահմանման՝ գոյություն ունի $U_i(x_0)\in \beta_{x_0}$, որ $U_i(x_0)\subset U$։ Այժմ ստանում ենք՝ $f(U_i(x_0))\subset f(U)\subset V_j(y_0 )$։
\par Հիմա հակառակը․ դիցուք $V$-ն $y_0$ կետի որևէ շրջակայք է։ Գոյություն ունի $V_j(y_0)\in \beta_{y_0}$, որ $V_j(y_0)\subset V$։ Ըստ պայմանի, գոյություն ունի $U_i(x_0)\in \beta_{x_0}$, որ $f(U_i(x_0))\subset V_j(y_0)\subset V$։ Տվյալ $V$ շրջակայքի համար որպես $U$ շրջակայք վերցնելով $U_i(x_0)$-ն ստանում ենք $f(U)\subset V$։ Ուստի $f$-ը անընդհատ է $x_0$ կետում։ 
\end{proof}
\par Կիրառենք այս թեորեմը մասնավոր դեպքում, երբ տոպոլոգիական տարածու\-թյունները մետրիկային են՝ $(X,\rho)$ և $(Y,\rho')$։
\par Ինչպես գիտենք, այդ տարածություններում բոլոր $\left\{\mathcal{D}(x,r)\right\}$ և $\left\{\mathcal{D}(y,r)\right\}$ անեզր գնդերի ընտանիքները կազմում են բազա համապատասխանաբար $x$ և $y$ կետերի շրջակայքերի համար (տե՛ս \hyperref[ch7thm2]{թեորեմ $2$}-ը թեմա $7$-ում)։ Այստեղից և թեորեմ $1$-ից ստանում ենք․

\begin{hetevanq}\label{ch9cor1} 
Մետրիկային տարածությունների $f:X \rightarrow Y$ արտապատկերու- \linebreak մը անընդհատ է ${x_0\in X}$ կետում այն և միայն դեպքում, երբ $Y$-ում ամեն մի \linebreak $\mathcal{D}(y_0,{r'})$ գնդի համար, որտեղ ${y_0=f(x_0)}$, գոյություն ունի $X$-ում $\mathcal{D}(x_0,r)$ գունդ, որ $f(\mathcal{D}(x_0,r))\subset \mathcal{D}(y_0,{r'})$։ 
\end{hetevanq}

\begin{definition}
Ունենք $X$ և $Y$ տոպոլոգիական տարածություններ և $f:X  \rightarrow  Y$ արտա\-պատկերում։ Ապա $f$-ը կոչվում է \textbf{անընդհատ արտապատկերում}, եթե այն անընդ\-հատ է $X$-ի բոլոր կետերում։
\end{definition}
Սկզբի համար բերենք անընդհատ արտապատկերումների երկու պարզ օրինակ։
\par ա) Ցանկացած տոպոլոգիական տարածության նույնական արտապատկերումը ինքն իր վրա անընդհատ արտապատկերում է։ Իրոք, ${f=\nuynakan_X:X  \rightarrow  X},\ f(x)=x$ արտա\-պատ\-կեր\-ման դեպքում $y_0=f(x_0)=x_0$ կետի $V$ շրջակայքի համար որպես $x_0$ կետի $U$ շրջակայք վերցնելով $V$-ն, կստանանք $f(U)=V\subset V$։ 
\par բ) Դիցուք $(X,\tau)$-ն և $(Y,\sigma )$-ն կամայական տոպոլոգիական տարածություններ են, $y_0\in Y$ սևեռված կետ է։ Ընդունված է $c:X  \rightarrow Y,\ c(x)=y_0,\ \forall x\in X$ արտա\-պատ\-կե\-րումն անվանել $X$-ի \textbf{հաստատուն արտապատկերում} $Y$-ի $y_0$ կետի վրա։ Այն անընդհատ արտապատ\-կերում է։ Իրոք, քանի որ $c^{-1}(y_0)=X$, ուստի $y_0$-ի $\forall\, V$ շրջակայքի համար որպես $ \forall x_0\in X$ կետի $U$ շրջակայք վերցնելով $X$-ը, կունենանք $f(U)\subset V$։

\begin{theorem}\label{ch9thm2} 
\label{թեորեմ $2$}
Ունենք $(X, \rho)$ և $(Y, \rho')$ մետրիկային տարածություններ։ Ապա \linebreak ${f:X  \rightarrow  Y}$ արտապատկերումն անընդհատ է այն և միայն այն դեպքում, երբ ցան\-կացած $x_0\in X$ կետի և $\varepsilon >0$ թվի համար գոյություն ունի $\delta=\delta(x_0, \varepsilon)>0$ թիվ, որ եթե $\rho(x,x_0)<\delta$, ապա ${\rho'} (f(x),f(x_0 ))< \varepsilon$։
\end{theorem}
Ապացուցումը անմիջականորեն հետևում է անընդհատության սահմանումից և \hyperref[ch9thm1]{թեորեմ $1$}-ի հետևանքից։
\begin{theorem}\label{ch9thm3}\label{թեորեմ $3$}
Ունենք $X$ և $Y$ տոպոլոգիական տարածություններ և ${f:X  \rightarrow  Y}$ արտապատկերում։ Ապա $f$-ը անընդհատ է այն և միայն այն դեպքում, երբ $Y\textrm{-ում}$ բաց ցանկացած ենթաբազմության նախակերպարը բաց ենթաբազմություն է $X\textrm{-ում}$։
\end{theorem}
\begin{proof}
\par ա) Ցույց տանք պայմանի անհրաժեշտությունը։ Դիտարկենք որևէ $x_0\in f^{-1} (V)$ կետ։ Քանի որ $f(x_0 )\in V$, ուստի $x_0$ կետում $f$-ի անընդ\-հա\-տութ\-յունից հետևում է, որ գոյություն ունի $x_0$-ի $U$ շրջակայք, որ $f(U)\subset V$։
\par Ունենք՝ $U\subset f^{-1}(f(U))\subset f^{-1} (V)$, որից հետևում է, որ $f^{-1} (V)$-ի ցանկացած $x_0$ կետ ներքին կետ է նրա համար, ուստի $f^{-1} (V)$-ն բաց ենթաբազմություն է $X$-ում։
\par բ) Այժմ ապացուցենք պայմանի բավարարությունը։ Դրա համար ցույց տանք, որ $f$-ը անընդհատ է ցանկացած $x_0\in X$ կետում։ Դիցուք $V$-ն $y_0=f(x_0)$ կետի որևէ շրջակայք է։ Գոյություն ունի $y_0$-ի $W$ բաց շրջակայք, որ $W\subset V$։ Ըստ պայմանի $f^{-1}(W)$-ն բաց ենթաբազմություն է $X$-ում։ Ուրեմն $U=f^{-1}(W)$-ն (բաց) շրջակայք է $x_0$-ի համար։ Ունենք՝ $f(U)=f(f^{-1}(W))\subset W\subset V$, ուստի $f$-ը անընդհատ է $x_0$ կետում։\qedhere
\end{proof}

\par Նշենք, որ թեորեմն ապացուցելիս մենք օգտվեցինք հետևյալ ներդրումներից (տե՛ս \hyperref[ch1thm3]{թեորեմ $3$}-ը թեմա $1$-ում)․ եթե ունենք $f:X  \rightarrow  Y$ արտապատկերում, $A\subset X$, $B\subset Y$ ենթաբազմություններ, ապա $f(f^{-1}(B))\subset B$,\, $A\subset f^{-1} (f(A))$։


\begin{hetevanq_counter}\label{ch9cor2}
Եթե հայտնի է $Y$ տարածության տոպոլոգիայի որևէ բազա, ապա ${f: X \to Y}$ արտապատկերման անընդհատության համար բավական է պահանջել, որ $X$-ում բաց լինեն տվյալ բազայի բոլոր տարրերի նախակերպարները (հիմնավո\-րե՛ք):
\end{hetevanq_counter}

\begin{hetevanq_counter}\label{ch9cor3} 
Եթե $f:X  \rightarrow  Y$, $g:Y  \rightarrow  Z$ արտապատկերումները անընդհատ են, ապա նրանց $g\circ f$ համադրույթը նույնպես անընդհատ է (\textbf{հիմնավորե՛ք})։
\end{hetevanq_counter}

\begin{theorem}[(անընդհատության հայտանիշ փակ ենթաբազմությունների տեր\-մին- \linebreak ներով)]\label{ch9thm4} \label{թեորեմ $4$}
Ունենք $X$ և $Y$ տոպոլոգիական տարածություններ։ Ապա որևէ ${f:X  \rightarrow  Y}$ արտապատկե\-րում անընդհատ է այն և միայն այն դեպքում, երբ $Y$-ում փակ ցան\-կա\-ցած ենթա\-բազ\-մութ\-յան նախակերպարը փակ ենթաբազմություն է $X$-ում։
\end{theorem}

\par Ապացուցելիս կօգտվենք ${f^{-1} (Y \setminus Z)=X \setminus f^{-1}} (Z)$ նույնությունից, որտեղ ${Z\subset Y}$։
\par ա) Ենթադրելով, որ $f$-ը անընդհատ է, կունենանք․ եթե $F$-ը փակ է $Y$-ում, ապա $(Y\setminus F)$-ը բաց է $Y$-ում, և ուրեմն $f^{-1} (Y  \setminus F)$-ը բաց է $X$-ում՝ ըստ \hyperref[ch9thm3]{թեորեմ 3}-ի։ Ուստի $f^{-1}(F)=X  \setminus f^{-1}(Y  \setminus F)$ ենթաբազմությունը փակ է $X$-ում։
\par բ) Հակառակը․ եթե $V$-ն բաց է $Y$-ում $ \Rightarrow  (Y  \setminus V)$-ն փակ է $Y$-ում $  \Rightarrow  f^{-1} (Y  \setminus V)$-ն փակ է $X$-ում $  \Rightarrow  f^{-1} (V)=X  \setminus f^{-1} (Y  \setminus V)$ ենթաբազմությունը բաց է $X$-ում։ Ուստի $f$-ը անընդհատ է ըստ \hyperref[թեորեմ $3$]{թեորեմ $3$}-ի։ \qed

\par Այժմ բերենք արտապատկերումների անընդհատության մի քանի հայտանիշ \linebreak ենթաբազմությունների փակման գործողության տերմիններով։
\par Ենթագիտակցաբար, $f:X \rightarrow Y$ անընդհատ արտապատկերումը մենք պատկե\-րացնում ենք որպես այնպիսի արտապատկերում, որը $X$-ի ցանկացած երկու «բա\-վա\-կա\-նաչափ մոտիկ» ($X$-ի տոպոլոգիայի իմաստով) $x_1$ և $x_2$ կետերի համա\-պա\-տաս\-խա\-նեցնում է $Y$-ի «բա\-վականաչափ մոտիկ» $f(x_1)$ և $f(x_2)$ կետեր ($Y$-ի տո\-պո\-լոգիայի իմաստով)։ Մյուս կողմից, ենթաբազմության հպման կետը մենք պատ\-կե\-րաց\-նում ենք որպես այդ ենթաբազմությանը «շատ մոտիկ», կամ նրան կպած կետ այն իմաստով, որ նրան հնարավոր չէ անջատել ենթաբազմությունից մի որևէ բաց շրջակայքով։ Ուստի ճշմարիտ է թվում հետևյալ պնդումը․ եթե $x \in X$ կետը հպման կետ է $A\subset X$ ենթաբազմության համար, ապա $f(x)$ կետը հպման կետ է $f(A)\subset Y$ ենթաբազմության համար։

\begin{theorem}[(անընդհատության հայտանիշ փակման գործողության տերմիններով)]\label{ch9thm5} \label{թեորեմ $5$}
\par $f:X \rightarrow  Y$ արտապատկերումը անընդհատ է այն և միայն այն դեպքում, երբ ցանկացած $A\subset X$ ենթաբազմության դեպքում $f( \widebar{A} )\subset \widebar{f(A)}$։
\end{theorem}

\begin{proof}
 \textbf{Պայմանի անհրաժեշտությունը։} Դիցուք $f$-ը անընդհատ է։ Վերց\-նենք կամայական $y\in f(\widebar{A})$ կետ և ցույց տանք, որ $y\in \widebar{f(A)}$։ Իրոք, գոյություն ունի $x\in \widebar{A}$ կետ, որ $f(x)=y$։ Դիտարկենք $y$ կետի կամայական $V$ շրջակայք։ Ըստ $x$ կետում $f$-ի անընդհատության՝ գոյություն ունի $x$-ի $U$ շրջակայք, որ $f(U)\subset V$։ Քանի որ $U\cap A\neq \varnothing $, ուստի $f(U)\cap f(A)\neq \varnothing \Rightarrow V  \cap f(A)\neq  \varnothing \Rightarrow y\in  \widebar{f(A)}$։ Այսպիսով՝ $f(\bar{A}) \subset  \widebar{f(A)}$։ 
 \par
 \textbf{Պայմանի բավարարությունը։} Դիտարկենք կամայական $B\subset Y$ փակ ենթա\-բազ\-մություն՝ $B=\widebar{B}$ և ապացուցենք, որ $A=f^{-1}(B)$  ենթաբազմությունը փակ է $X$-ում (դրանից կհետևի, որ $f$-ը անընդհատ է ըստ \hyperref[թեորեմ $4$]{թեորեմ $4$}-ի)։

\par Դրա համար ցույց տանք, որ $A$-ն պարունակում է իր բոլոր հպման կետերը։ Եթե $x\in \widebar{A}$, ապա ըստ պայմանի՝ $f(x)\in f(\widebar{A})\subset \widebar{f(A)}\subset \widebar{B}=B$: Ուստի $x \in {f^{-1}(B)}=A$, որից էլ ստանում ենք $\widebar{A}=A$։ Ուրեմն $A$-ն փակ է $X$-ում։
\end{proof}

\par Գոյություն ունի նաև արտապատկերման անընդհատության հայտանիշ ենթա\-բազ\-մության ներքինի տերմիններով։
\begin{theorem}\label{ch9thm6} 
${f:X  \rightarrow Y}$ արտապատկերումը անընդհատ է այն և միայն այն դեպ\-քում, երբ կամայական $B\subset Y$ ենթաբազմության համար $f^{-1} (\inter B)\subset \inter (f^{-1} (B))$։
\end{theorem}

\begin{proof}
\par ա) Ենթադրենք $f$-ը անընդհատ է։ Վերցնենք կամայական $x\in f^{-1} (\inter B)$ կետ և ցույց տանք, որ $x\in \inter(f^{-1} (B))$։ Ունենք՝ $f(x)\in \inter B\subset B$, ուրեմն $x\in f^{-1} (\inter B)\subset f^{-1} (B)$։ Քանի որ, ըստ \hyperref[ch9thm3]{թեորեմ $3$}-ի, $f^{-1} (\inter B)$-ն բաց ենթաբազմություն է $X$-ում, ուստի $x$ կետը ներքին կետ է $f^{-1} (B)$-ի համար  $\Rightarrow$ \linebreak  ${x\in \inter(f^{-1} (B))}$։
\par բ) Հակառակ պնդումը ապացուցելու համար վերցնենք $\forall  V\subset Y$ բաց ենթա\-բազ\-մություն և ցույց տանք, որ $f^{-1}$ $(V)$-ն բաց է $X$-ում։ Շարունակությունը թողնում ենք ընթերցողին։
\end{proof}

\begin{definition}
Արտապատկերումը կոչվում է \textbf{խզվող}, եթե այն անընդհատ չէ։
\end{definition}

\par Բերենք խզվող արտապատկերումների օրինակներ՝ հիմնավորումները կատա\-րե\-լով \hyperref[ch9thm5]{թեորեմ $5$}-ում բերված հայտանիշի միջոցով։

\begin{example}\label{ch9ex1}
Դիտարկենք $(\R, \textrm{սովոր․})$ տարածությունը և ${f:\R \rightarrow \R}$ արտա\-պատ\-կե\-րումը, որտեղ $f(x)$-ը $x\in \R$ թվի ամբողջ մասն է։ Ըստ սահմանման՝ եթե \linebreak ${x\in [n, n+1)},\ {n\in \Z}$, ապա ${f(x)=n}$։ Մասնավորապես ${f(n)=n},\ \forall {n \in \Z}$։ Վերց\-նե\-լով $A=[n-1, n)$՝ կունենանք $f(A)=\{n-1\}$ և $\widebar A =[n-1, n]$։ Բացի այդ $n\in \widebar{A} \Rightarrow f(n)\in f(\widebar{A})$։ Ենթադրելով, որ $f$-ը անընդհատ է $n\in Z$ կետում, կստա\-նանք՝ $f(n)\in f(\widebar{A})\subset \widebar{f(A)} =\widebar{\{n-1\}}=\{n-1\}$։ Ստացանք $f(n)=n-1$ (հակա\-սու\-թյուն)։ Հետևաբար $f$-ը խզվող է, անընդհատ չէ $\R$-ի $n\in \Z$ կետերում։
\end{example}
\begin{example}\label{ch9ex2} 
Մաթեմատիկական անալիզի դասընթացում բերվում է թվային ֆունկ\-ցիայի օրինակ (Դիրիխլեի ֆունկցիան), որն անընդհատ չէ $\R$-ի բոլոր կետերում։ Այն սահմանվում է որպես $f:\R \rightarrow \R$ արտապատկերում, որտեղ $f(x)=0$, երբ $x$-ը ռացիոնալ է՝ $x\in \Q$, և $f(x)=1$, երբ $x$-ը իռացիոնալ է՝ $x\in {\rm Ir}$։ Նկատենք, որ $\R=\Q\cup {\rm Ir}$, $\Q\cap {\rm Ir}=\varnothing$, $\widebar{\Q}=\widebar{{\rm Ir}}=\R$։ Այդ արտապատկերումը ընդհանրացվում է հետևյալ տեսքով։ 

\par Դիցուք $X$-ը տոպոլոգիական տարածություն է, ընդ որում գոյություն ունեն $X$-ի այնպիսի $A$ և $B$ ենթաբազմություններ, որ $X=A\cup B$, $A\cap B=\varnothing$, $\widebar{A}=\widebar{B}=X$։ Ապա $f:X \rightarrow \R$ արտապատկերումը, որտեղ $f(x)=0$, երբ $x\in A$ և $f(x)=1$, երբ $x\in B$, անընդհատ չէ $X$-ի բոլոր կետերում։ Ապացուցելու համար վերցնենք որևէ $x_0\in A$ կետ։ Ունենք $f(x_0)=0$։ Մյուս կողմից, քանի որ $X=\widebar{B}$, ստանում ենք՝ $x_0\in \widebar{B}$, որից հետևում է $f(x_0)\in f(\widebar{B})\subset  \widebar{f(B)}=\widebar{\{1\}}=\{1\}$։ Ստացանք $f(x_0)=1$ (հակասություն)։ Նման ձևով ապացուցվում է, որ $f$-ը անընդհատ չէ նաև $B$-ի բոլոր կետերում։
\end{example}

\par Այժմ պատրաստ ենք ներմուծել մի կարևորագույն հասկացություն։ Նկատենք, որ ցանկացած հանրահաշվական կամ երկրաչափական տեսություն կառուցելիս, որի նպատակը ուսումնասիրվող օբյեկտների (լինեն դրանք խմբեր, օղակներ, գծ․ տարածություններ, երկրաչափական պատկերներ և այլն) դասակարգումն է, հստա\-կեց\-վում է հետևյալ հարցը․ ե՞րբ են տվյալ տեսության մեջ ուսումնասիրվող երկու օբյեկտներ համարվում նույնը կամ նույնականացվող։ Օրինակ՝ խմբերի տեսութ\-յու\-նում երկու $G_1$ և $G_2$ խմբեր համարվում են նույնը, եթե նրանք իզոմորֆ են։ Դա համարժեք է նրան, որ գոյություն ունենան ${\varphi_1:G_1\rightarrow G_2}$ և $\varphi_2:G_2\rightarrow G_1$ հոմոմորֆիզմներ այնպես, որ $\varphi_2\circ \varphi_1=\nuynakan_{G_1}$ և $\varphi_1\circ \varphi_2=\nuynakan_{G_2}$։ Տոպոլոգիայում հոմոմորֆիզմների դերը կատարում են հոմեոմորֆիզմները։ Այդ երկու անվանում\-ները գրեթե համահնչուն են, բայց ունեն միանգամայն տարբեր իմաստ և գործածե\-լիս պետք է չշփոթել։

\par Ինչպես գիտենք (տե՛ս \hyperlink{sec:1}{թեմա 1}-ում), ամեն մի փոխմիարժեք $f:X\rightarrow Y$ արտա\-պատ\-կեր\-ման համար գոյություն ունի նրան հակադարձ $f^{-1}:Y\rightarrow X$ ար\-տա\-պատ\-կե\-րում, ընդ որում $f\circ f^{-1}=\nuynakan_Y$ և $f^{-1}\circ f=\nuynakan_X$։ Այս դեպքում ասում են նաև, որ $f$ արտապատկերումը հակադարձելի է։ Բերենք հակադարձելիությանը համարժեք պայման (հայտանիշ)։ 

\par Որևէ $f:X\rightarrow Y$ արտապատկերում հակադարձելի է այն և միայն այն դեպքում, երբ գոյություն ունի $g:Y\rightarrow X$ արտապատկերում, որ $f\circ g=\nuynakan_Y$ և $g\circ f=\nuynakan_X$։
\par Իրոք, եթե $f$-ը հակադարձելի է, ապա որպես $g:Y\rightarrow X$ արտապատկերում կարելի է վերցնել $f^{-1}$-ը։ Հակառակը․ դիցուք $f$-ի համար $\exists g:Y\rightarrow X$, որ $f\circ g=\nuynakan_Y$ և $g\circ f=\nuynakan_X$։ Ցույց տանք, որ $f$-ը փոխմիարժեք է։ Ենթադրենք $f(x_1)=f(x_2)$։ Ունենք՝ $x_1=\nuynakan_X(x_1)=g\circ f(x_1)=g(f(x_1))$, նման ձևով $x_2=g(f(x_2))$, ուստի $x_1=x_2$, այսինքն $f$-ը ինյեկտիվ է։ Կամայական $y\in Y$ կետի համար ունենք՝ $y=\nuynakan_Y(y)=f\circ g(y)=f(g(y))$։ Նշանակելով $g(y)=x\in X$, ստանում ենք $f(x)=y$, այսինքն $f$-ը սյուրյեկտիվ է։ Ուստի $f$-ը փոխմիարժեք արտապատկերում է։ \qed

% բերված 11-րդ թեմայից - դեռ չխմբագրած


\begin{definition}
Ունենք $X$, $Y$ տոպոլոգիական տարածություններ։ Ապա $f:X\rightarrow Y$ հա\-կա\-դար\-ձելի արտապատկերումը կոչվում է \textbf{հոմեոմորֆիզմ}, եթե $f$-ը և նրա հա\-կա\-դարձ $f^{-1}:X\rightarrow Y$ արտապատկերումը անընդհատ են։ Ասում են նաև, որ $X$ տարածությունը \textbf{հոմեոմորֆ} է $Y$ տարածությանը, եթե գոյություն ունի որևէ ${f:X\rightarrow Y}$ հոմեոմորֆիզմ։ 
\end{definition}
Սահմանումից հետևում է․
\begin{enumerate}
    \item [1)] Ցանկացած $X$ տոպոլոգիական տարածության համար $\nuynakan_X:X\rightarrow X$ նույնական արտապատկերումը հոմեոմորֆիզմ է։ Հետևաբար ամեն մի տոպոլոգիական տարածություն հոմեոմորֆ է ինքն իրեն։
    \item [2)] Եթե $f$-ը հոմեոմորֆիզմ է, ապա $f^{-1}$-ը ևս հոմեոմորֆիզմ է։\\Ուստի, եթե $X$-ը հոմեոմորֆ է $Y$-ին, ապա իր հերթին $Y$-ը հոմեոմորֆ է $X$-ին։
\end{enumerate}
\par Վերը շարադրվածից նաև հետևում է․ $X$ և $Y$ տոպոլոգիական տարածություննե\-րը հոմեոմորֆ են այն և միայն այն դեպքում, երբ գոյություն ունեն $f:X\rightarrow Y$ և $g:Y\rightarrow X$ անընդհատ արտապատկերումներ, որ $f\circ g=\nuynakan_Y$ և $g\circ f=\nuynakan_X$։
\begin{theorem}\label{ch9thm7} 
Տոպոլոգիական տարածությունների հոմեոմորֆության հա\-րա\-բե\-րու\-թյունը համարժեքության հարաբերություն է։ 
\end{theorem}

\begin{varproof}
Ապացուցման համար մնացել է ստուգել տրան\-զի\-տի\-վութ\-յան (փոխանցականու\-թյան) հատկությունը։ Դիցուք $X$-ը հոմեոմորֆ է $Y$-ին, և $Y$-ը հոմեոմորֆ է $Z$-ին։ Նշանակում է գոյություն ունեն $f:X\rightarrow Y$ և $g:Y\rightarrow Z$ հոմեոմորֆիզմներ։ Պարզ է, որ $g\circ f:X\rightarrow Z$ արտապատկերումը փոխմիարժեք է որպես փոխմիարժեք ար\-տա\-պատկերումների համադրույթ։ Ուստի այն հակադարձելի է։ Այն նաև անընդ\-հատ է որպես անընդհատ արտապատկերում\-ների համադրույթ։ Բացի այդ $(g\circ f)^{-1}=f^{-1}\circ g^{-1}$ արտապատկերումը $Z$-ից $X$ անընդհատ է նույն հիմքով։ Հետևաբար $g\circ f$ արտապատկերումը հոմեոմորֆիզմ է։ Ուստի $X$-ը հոմեոմորֆ է $Z$-ին։ \end{varproof}

\par Տոպոլոգիական տարածությունների համարժեքության (ըստ հոմեոմորֆության) ամեն մի դաս կոչվում է \textbf{տոպոլոգիական տիպ}։ Ընդհանուր տոպոլոգիայում տվյալ տոպոլոգիական տիպը ներկայացնող բոլոր տարածությունները (միմյանց հոմեո\-մորֆ տարածությունները) համարվում են միատեսակ (նույնը կամ նույնականաց\-վող)։

\begin{definition} 
Տոպոլոգիական տարածության որևէ հատկություն կոչվում է \textbf{տո\-պո\-լո\-գիական հատկություն} կամ \textbf{տոպոլոգիական ինվարիանտ}, եթե այդ հատ\-կութ\-յամբ օժտված են միմյանց հոմեոմորֆ տվյալ պահին դիտարկվող բոլոր տո\-պո\-լո\-գիական տարածությունները։
\end{definition}

% \par Այսպիսով, յուրաքանչյուր տոպոլոգիական տիպ կարելի է ներկայացնել որպես այն բոլոր տոպոլոգիական հատկությունների համախմբությունը, որոնցով օժտված են տվյալ պահին դիտարկվող այդ տիպին պատկանող բոլոր տարածությունները։

\par Այժմ հարցին, թե ի վերջո ինչ է տոպոլոգիան, կարելի է պատասխանել այսպես։ Տոպոլոգիան գիտություն է տոպոլոգիական տարածությունների (երկրաչափական պատ\-կերների) և այդ տարածությունների անընդհատ արտապատկերումների (երկ\-րա\-չափական պատկերների անընդհատ ձևափոխությունների) տոպոլոգիական \linebreak հատ\-կությունների մասին։ 

\par Վերոհիշյալից հետևում է, որ երկու տոպոլոգիական տարածություն հոմեոմորֆ չեն այն և միայն այն դեպքում, երբ գոյություն ունի տոպոլոգիական հատկություն, որով օժտված է այդ տարածություններից մեկը, բայց օժտված չէ մյուսը։
\par Ընդհանուր տոպոլոգիայի հիմնական խնդիրը կարելի է կարճ ձևակերպել այս\-պես․ կառուցել տոպոլոգիական ինվարիանտների լրիվ համակարգ տվյալ պահին դի\-տարկ\-վող բոլոր տոպոլոգիական տարածությունների համար։
\par Տոպոլոգիական հատկությունների օրինակներ են անջատելիության աքսիոմնե\-րը, հաշվելիության առաջին և երկրորդ աքսիոմները, սեպարաբելությունը և այլն (հիմնավորե՛ք)։

\par Հոմեոմորֆ (ոչ հոմեոմորֆ) տարածությունների երկրաչափական բնույթի օրի\-նակներ, ինչպես նաև տո\-պո\-լո\-գիական այլ ինվարիանտների օրինակներ կբերենք խնդիրների բաժնում և հաջորդ թեմաներում։ Իսկ հիմա, որպես օրինակ ցույց տանք, որ սեպարաբելությունը տոպոլոգիական հատկու\-թյուն է։ Այդ նպատակով նախ ա\-պա\-ցուցենք՝
\begin{theorem}\label{ch9thm8}
\label{11:թեորեմ 4}
    Եթե $X$ և $Y$ տարածություններից $X$-ը սեպարաբել է և գոյություն ունի անընդհատ, սյուրյեկտիվ $f:X\to Y$ արտապատկերում, ապա $Y$-ը ևս սե\-պա\-րաբել է։
\end{theorem}
\begin{proof}
Ըստ պայմանի՝ $X$-ում գոյություն ունի հաշվելի, ամենուրեք խիտ $\{x_n\}$ ենթաբազմություն։ Ապացուցենք, որ $\{f(x_n)\}$ հաշվելի ենթաբազմությունը ա\-մենուրեք խիտ է $Y$-ում։ Դիտարկենք կամայական $V\subset Y$ բաց ենթաբազմություն և ցույց տանք, որ $\{f(x_n)\} \cap V \neq \varnothing$ (դրանից կհետևի $Y$-ի սեպարաբելությունը ըստ թեմա 8-ում \hyperref[ch8thm4]{թեորեմ 4}-ի)։ Քանի որ $f$-ը անընդհատ է, ուստի $f^{-1}(V)$-ն բաց են\-թաբազմություն է $X$-ում։ Այժմ $X$-ի սեպարաբելությունից հետևում է $\{x_n\}\cap f^{-1}(V)\neq \varnothing$, և ուրեմն նաև $f(\{x_n\})\cap f(f^{-1}(V))\neq \varnothing$։ Քանի որ $f(f^{-1}(V))=V$ (հետևում է $f$-ի սյուրյեկտիվությունից՝ $f(X)=Y$), ուստի $f(\{x_n\})\cap V\neq \varnothing$։
\end{proof}
\begin{hetevanq8_sa}\label{ch11cor1}
Եթե միմյանց հոմեոմորֆ երկու տարածություններից մեկը սեպարաբել է, ապա մյուսը նույնպես սեպարաբել է։ Իրոք, դիցուք $X$-ը սե\-պա\-րա\-բել տարածություն է և հոմեոմորֆ է $Y$-ին։ Ուստի գոյություն ունի $f:{X\to Y}$ հոմեոմորֆիզմ։ Հոմեոմորֆիզմի վերը բերված սահմանումից մասնավորապես հե\-տևում է, որ $f$-ը անընդհատ, սյուրյեկտիվ արտապատկերում է, և ուրեմն $Y$-ը սեպա\-րաբել է համաձայն \hyperref[11:թեորեմ 4]{թեորեմ 8}-ի։
\end{hetevanq8_sa}

% \subsection*{Բաց և փակ արտապատկերումներ}
\par Ինչպես գիտենք, եթե $f:X\rightarrow Y$ արտապատկերման դեպքում $Y$-ում բաց բոլոր ենթաբազմությունների նախակերպարները բաց ենթաբազմություններ են $X$-ում, ապա $f$-ը անընդհատ է։
\par Իսկ ի՞նչ կարելի է ասել այն արտապատկերումների մասին, որոնք հակառակը՝ $X$-ում բաց ենթաբազմությունները արտապատկերում են $Y$-ում բաց են\-թա\-բազ\-մութ\-յունների։ Հետևյալ հասարակ օրինակը ցույց է տալիս, որ այդպիսի ար\-տա\-պատ\-կե\-րումը կարող է  անընդհատ չլինել։ Դիտարկենք $f:(X, \textrm{անտ․}) \rightarrow (X, \textrm{դիսկր․})$ արտապատկերումը, որտեղ $X$-ը պարունակում է մեկից ավելի կետեր, իսկ $f$-ը $X$-ի նույնական արտապատկերումն է։ Պարզ է, որ $f$-ը անընդհատ չէ (ինչո՞ւ), չնայած որ այն $X$-ում բաց ենթաբազմությունները արտապատկերում է $Y$-ում բաց ենթա\-բազ\-մու\-թյունների։ Այնուամենայնիվ, նշված հատկության և անընդհատության համակ\-ցումը բերում է բաց (փակ) արտապատկերումներ հասկացություններին, որոնք օգ\-տակար են և հարմար որոշ իրավիճակներում։

\begin{definition} 
$f:X\rightarrow Y$ անընդհատ արտապատկերումը կոչվում է \textbf{բաց ար\-տա\-պատ\-կերում}, եթե $X$-ում բաց ցանկացած ենթաբազմության կերպարը բաց ենթա\-բազ\-մու\-թյուն է $Y$-ում։
\par Նման ձևով սահմանվում է \textbf{փակ արտապատկերումը}՝ նախորդ սահմանման \linebreak մեջ ամենուրեք «բաց» բառը փոխարինելով «փակ» բառով։
\end{definition}
\par Նկատենք, որ վերը բերված օրինակում $\nuynakan_X:(X, \textrm{անտիդ․})\rightarrow (X, \textrm{դիսկր․})$-ը և՛ բաց, և՛ փակ արտապատկերում է, իսկ $\nuynakan_X:(X, \textrm{դիսկր․})\rightarrow (X,\textrm{անտիդ․})$ արտա\-պատ\-կերումը ոչ բաց և ոչ էլ փակ արտապատկերում է։

\begin{example}\label{ch9ex3}
Հետևյալ՝ $c:(\R,\textrm{սովոր․}) \rightarrow (\R, \textrm{սովոր․}),\ c(x)=0,\ \forall x\in \R$ հաստա\-տուն արտապատկերումը $0$ կետի վրա, փակ արտապատկերում է, բայց բաց ար\-տա\-պատ\-կե\-րում չէ (հետևում է նրանից, որ $\{0\}\subset \R$ ենթաբազմությունը փակ, բայց ոչ բաց ենթաբազմություն է)։
\end{example}

\begin{example}\label{ch9ex4} 
Ցույց տանք, որ $\R^2$ կոորդինատային հարթության $P:\R^2\rightarrow \R$, $P(x_1,x_2)=x_1$ պրոյե\-կցիան բաց, բայց ոչ փակ արտապատկերում է։ Իրոք, $\R^2\textrm{-ի}$ մետրիկային տոպոլոգիայի համար բազա են կազմում բոլոր անեզր շրջանները, իսկ $(\R, \textrm{սովոր․})$-ի համար՝ բոլոր անեզր $(a,b)$ ինտերվալները։

\par Քանի որ անեզր շրջանի պրոյեկցիան անեզր ինտերվալ է, ուստի $P$-ն $\R^2$-ում բաց ենթաբազմությունները արտապատկերում է $\R$-ում բաց ենթաբազմությունների։ Բա\-ցի այդ, $P$-ն անընդհատ է, քանի որ $P^{-1} (a,b)=\R \times(a,b)$ ենթաբազմությունները բաց են $\R^2$-ում։ Այսպիսով $P$-ն բաց արտապատկերում է։ Ցույց տանք, որ $P$-ն փակ արտապատկերում չէ։ 

\par Դիտարկենք $A=\left\{\left(x;\, \dfrac{1}{x}\right)\mid x\in\R\setminus\{0\}\right\}$ ենթաբազմությունը $\R^2$-ում ($y=\dfrac{1}{x}$ հիպերբոլի գրաֆիկը)։ Այն փակ ենթաբազմություն է, քանի որ $\R^2\setminus A$ լրացումը բաց է (հիմնավորե՛ք)։ Բայց նրա $P(A)=(-\infty,0) \cup (0, +\infty)$ կերպարը փակ չէ $\R$-ում (\textbf{ինչո՞ւ})։
Ուստի՝ $P$-ն փակ արտապատկերում չէ։ \qed

\begin{center}
\begin{tikzpicture}
\draw[->, ultra thick, blue!70] (0,3) -- (3,3);
\draw[->, ultra thick, blue!70] (6,3) -- (3,3);
\draw[->,  thick] (3,0) -- (3,6);

\node[anchor=east] at (0, 3) {$\R$};
\node[anchor=north east] at (2, 5) {$\R^2$};
\node[anchor=north west] at (3.05,2.95) {$0$};

\draw[scale=1, domain=3.4:5.6, smooth, variable=\x, blue] plot ({\x}, {1/(\x-3)+3});
\draw[scale=1, domain=0.4:2.6, smooth, variable=\x, blue] plot ({\x}, {1/(\x-3)+3});

\draw [-stealth](3.5,4) -- (3.5,3.25) node[anchor=south west] {$P$};
\draw [-stealth](2.5,2) -- (2.5,2.75) node[anchor=north east] {$P$};

\end{tikzpicture}
\end{center}

\end{example}
\par Ինչպես գիտենք, եթե $f:X\rightarrow Y$ և $g:Y\rightarrow Z$ արտապատկերումներն անընդհատ են, ապա նրանց $g\circ f:Y\rightarrow Z$ համադրույթը անընդհատ է։ Քննարկենք հակառակ խնդիրը․ դիցուք հայտնի է, որ $g\circ f$ համադրույթը անընդհատ է, նաև անընդհատ է $f$ կամ $g$ արտապատկերումներից մեկը։ Կլինի՞ անընդհատ մյուսը։ Հետևյալ օրի\-նակ\-նե\-րը ցույց են տալիս, որ ընդհանուր դեպքում հարցի պատասխանը բացա\-սա\-կան է։

\begin{example}\label{ch9ex5} 
Դիտարկենք որևէ $f:X\rightarrow Y$ ոչ անընդհատ արտապատկերում և ${g:Y\rightarrow \R},\ g(y)=0,\ y\in\R$ հաստատուն արտապատկերումը։ Պարզ է, որ ${g\circ f:X\rightarrow\R}$ համադրույթը նույնպես հաստատուն արտապատկերում է, և հետևա\-բար այն անընդհատ է։ Այսպիսով՝ $g\circ f$ –ը և $g$-ն անընդհատ են, բայց $f$-ը անընդհատ չէ։
\end{example}

\par Ընթերցողին առաջարկում ենք բերել նման օրինակ, երբ անընդհատ են $f:X\rightarrow Y$ և $f\circ g:X\rightarrow Z$ արտապատկերումները, բայց անընդհատ չէ $g:Y\rightarrow Z$ արտապատկերումը։
\par Հետևյալ երկու հատուկ դեպքերում $f$, $g$ արտապատկերումներից մեկի և նրանց $g\circ f$ համադրույթի անընդհատությունից հետևում է նաև մյուսի անընդհատությունը։

\begin{theorem}\label{ch9thm9}
Ունենք $f:X\rightarrow Y,\ g:Y\rightarrow Z$ արտապատկերումներ, ընդ որում $g\circ f$-ը անընդհատ է․
\begin{enumerate}
    \item [ա)] եթե $g$-ն ինյեկտիվ բաց (կամ փակ) արտապատկերում է, ապա $f$-ը անընդհատ է,
    \item [բ)] եթե $f$-ը սյուրյեկտիվ բաց (կամ փակ) արտապատկերում է, ապա $g$-ն անընդ\-հատ է։
\end{enumerate}
\end{theorem}

\begin{varproof}
Ապացուցենք ա)-ն։ Դիցուք $V\subset Y$ ենթաբազմությունը բաց է $Y$-ում, ցույց \linebreak տանք, որ $f^{-1}(V)$-ն բաց է $X$-ում։ Ըստ պայմանի $g(V)$-ն բաց է $Z$-ում։ Հետևաբար $(g\circ f)^{-1}(g(V))$-ն բաց է $X$-ում։ Ունենք՝ $(g\circ f)^{-1}(g(V))=f^{-1}\big(g^{-1}(g(V))\big)=f^{-1} (V)$ $\Rightarrow f^{-1} (V)$-ն բաց է $X\textrm{-ում} \Rightarrow f$-ը անընդհատ է։
\end{varproof}

\par Նման ձևով ապացուցվում է բ)-ն (ապացուցումը թողնում ենք ընթերցողին)։ 
\par Բերենք նաև հոմեոմորֆիզմի հայտանիշ բաց (փակ) արտապատկերումների տերմիններով։

\begin{theorem}\label{ch9thm10} 
Տոպոլոգիական տարածությունների որևէ $f:X\rightarrow Y$ արտապատ\-կե\-րում հոմեոմորֆիզմ է այն և միայն այն դեպքում, երբ $f$-ը բաց (կամ փակ) բիյեկտիվ արտապատկերում է։
\end{theorem}

\begin{proof} 
Դիցուք $f$-ը հոմեոմորֆիզմ է։ Նշանակում է $f$-ը անընդհատ է, փոխմիարժեք է և նրա $g:Y\rightarrow X$ հակադարձ արտապատկերումը նույնպես անընդ\-հատ է։ Եթե $U$-ն որևէ բաց ենթաբազմություն է $X$-ում, ապա $f$-ի բիյեկտիվությունից և $g$-ի անընդհատությունից հետևում է, որ $f(U)=g^{-1}(U)$ ենթաբազմությունը բաց է $Y\textrm{-ում}$, ուստի $f$-ը բաց արտապատկերում է։ $f$-ի բիյեկտիվությունից հետևում է, որ $f$-ը նաև փակ արտապատկերում է։ 

\par Այժմ հակառակը․ դիցուք $f$-ը բաց, բիյեկտիվ արտապատկերում է։ Նշանակում է $f$-ը անընդհատ է, և գոյություն ունի նրա հակադարձ $g:Y\rightarrow X$ արտապատկերում։ Մնում է ցույց տալ, որ $g$ արտապատկերումը անընդհատ է։ Եթե $U$-ն որևէ բաց են\-թա\-բազ\-մություն է $X$-ում, ապա $g^{-1}(U)=f(U)$ ենթաբազմությունը բաց է $Y$-ում ըստ պայմանի։
Ուստի $g$-ն անընդհատ է ըստ \hyperref[ch9thm3]{թեորեմ 3}-ի։
\end{proof}


\bigskip
\bigskip
\subsubsection*{Խնդիրներ և հարցեր թեմա 9-ի վերաբերյալ}

\begin{enumerate}[label=\thesection.\arabic*.]

% 9.1
\item Դիցուք ունենք կամայական $f: X \to Y$ անընդհատ արտապատկերում։ Ճի՞շտ է արդյոք պնդումը․ ամեն մի $x_0 \in X$ կետի ցանկացած $U$ շրջակայքի $f(U)$ կերպարը շրջակայք է $f(x_0) \in Y$ կետի համար։

\begin{hint}
Դիտարկեք որևէ $f: X \to Y$ հաստատուն արտապատկերում, որ\-տեղ $Y$-ը օժտված է անտիդիսկրետ տոպոլոգիայով և պարունակում է մեկից ավելի թվով տարրեր։
\end{hint}

% 9.2
\item Դիցուք $X$ տոպոլոգիական տարածությունը օժտված է հետևյալ հատկությամբ․ ցանկացած $Y$ տոպոլոգիական տարածության և ցանկացած $f: X \to Y$ արտա\-պատկերման դեպքում $f$-ը անընդհատ է։ Ապացուցեք, որ $X$-ի տոպոլո\-գիան դիսկրետն է։

\begin{hint}
Որպես $Y$ վերցրեք $X$ բազմությունը դիսկրետ տոպոլոգիայով։
\end{hint}

% 9.3
\item Դիցուք $Y$ տոպոլոգիական տարածությունը օժտված է հետևյալ հատկությամբ․ ցանկացած $X$ տոպոլոգիական տարածության և ցանկացած $f: X \to Y$ արտա\-պատկերման դեպքում $f$-ը անընդհատ է։ Ապացուցեք, որ $Y$-ի տոպոլո\-գիան անտիդիսկրետն է։

\begin{hint}
Որպես $X$ վերցրեք $Y$ բազմությունը անտիդիսկրետ տոպոլոգիա\-յով։
\end{hint}

% 9.4
\item Դիցուք $f: X \to Y$ արտապատկերումն այնպիսին է, որ $X$-ում բաց ցանկացած ենթաբազմության կերպարը բաց ենթաբազմություն է $Y$-ում։ Ապացուցեք, որ այդպիսի արտապատկերումը կարող է անընդհատ չլինել։

\begin{hint}
Դիտարկեք որևէ $X$ բազմություն $\tau$ և $\sigma$ տոպոլոգիաներով, որտեղ $\tau \subset \sigma$, իսկ որպես $f$ վերցրեք $\nuynakan_X$ նույնական արտապատկերումը։
\end{hint}

% 9.5
\item Բերեք $X$, $Y$ տոպոլոգիական տարածությունների և $f: X \to Y$ բիյեկտիվ ան\-ընդհատ արտապատկերման օրինակ, որ $f^{-1}$ հակադարձ արտապատկերումը չլինի անընդհատ։

\begin{hint}
Օգտվեք նախորդ խնդրի ցուցումից։
\end{hint}

% 9.6
\item Դիցուք $\R^2$ էվկլիդեսյան հարթության ինչ-որ $f: \R^2 \to \R^2$ արտապատ\-կե\-րում բավարարում է $\rho\left(f(x_1), f(x_2)\right) = r \cdot \rho(x_1, x_2)$ պայմանին $\forall \ x_1, x_2 \in \R^2$ կետերի դեպքում, որտեղ $r$-ը սևեռված իրական դրական թիվ է։ Ապա\-ցու\-ցեք, որ $f$-ը անընդհատ արտապատկերում է։

% 9.7
\item Ունենք $\R$ թվային ուղղի $g: \R \to \R$ անընդհատ արտապատկերում։ Հայտնի է, որ $g(x)$-ը ամբողջ թիվ է ցանկացած $x \in \R$ կետի դեպքում։ Ապացուցեք, որ $g$-ն հաստատուն արտապատկերում է։

% 9.8
\item Ունենք $f: X \to \R$ անընդհատ արտապատկերում, որտեղ $X$-ը կամայական տոպոլոգիական տարածություն է, և $a \in \R$ սևեռված թիվ։ Ապացուցեք, որ $g: X \to \R$ արտապատկերումը՝ սահմանված $g(x) = a f(x), \ x \in X$ բանաձևով, նույնպես անընդհատ է։

\begin{hint}
Ներկայացրեք $g$-ն որպես $f$ և $h: \mathbb{R} \to \mathbb{R}, \ h(r) = ar, \ r \in \mathbb{R}$ արտապատկերումների համադրույթ:
\end{hint}

% 9.9
\item Ունենք $X$ բազմության որևէ երկու $\tau$ և $\sigma$ տոպոլոգիա։ Ապացուցեք․ $X \to X$ նույնական արտապատկերումը $(X, \tau)$ և $(X, \sigma)$ տարածությունների հոմեոմոր\-ֆիզմ է այն և միայն այն դեպքում, երբ $\tau$-ն և $\sigma$-ն նույնն են։

% 9.10
\item Դիցուք $f: X \to Y$ արտապատկերումը $(X, \tau)$ և $(Y, \sigma)$ տոպոլոգիական տա\-րա\-ծությունների հոմեոմորֆիզմ է։ Ապացուցեք․ եթե որևէ $U \subset X$ ենթաբազմու\-թյուն շրջակայք է որևէ $x_0 \in X$ կետի համար, ապա $f(U) \subset Y$ ենթաբազմու\-թյունը շրջակայք է $f(x_0)$ կետի համար։

% 9.11
\item Ճի՞շտ է արդյոք պնդումը․ վերջավոր թվով տարրերից կազմված տոպոլոգիա\-կան տարածությունների դեպքում տարրերի քանակությունը տոպոլոգիական ինվարիանտ է։

% 9.12
\item Ընտրեք անջատելիության $\textrm{T}_0$, $\textrm{T}_1$, $\textrm{T}_2$ աքսիոմներից որևէ մեկը և ապացուցեք, որ այն տոպոլոգիական հատկություն է։

% 9.13
\item Բերեք տոպոլոգիական տարածությունների որևէ $f: X \to Y$ փակ, բայց ոչ բաց արտապատկերման օրինակ։

\begin{hint}
Դիտարկեք $\R $ թվային ուղիղը և $\R ^2$ կոորդինատային հարթու\-թյու\-նը սովորական (էվկլիդյան) մետրիկային տոպոլոգիաներով։ Ապացուցեք, որ $f: \R  \to \R ^2, \ f(x)=(x, 0), \ x \in \R $ արտապատկերումը փակ, բայց ոչ բաց արտապատկերում է։
\end{hint}

% 9.14
\item Ունենք $f: X \to Y$ և $g: Y \to Z$ արտապատկերումներ, ընդ որում ${g \circ f: X \to Z}$ համադրույթը անընդհատ է։ Ապացուցեք․ եթե $f$-ը սյուրյեկտիվ բաց (կամ փակ) արտապատկերում է, ապա $g$-ն անընդհատ արտապատկերում է։

% 9.15
\item Ճի՞շտ է արդյոք, որ թվային ուղղի $\R  \to \R $ նույնական արտապատկերումը որո\-շում է տոպոլոգիական տարածությունների $(\R, \mapsfrom) \to (\R, \mapsto)$ հոմեոմոր\-ֆիզմ, որտեղ $\mapsfrom$ և $\mapsto$ սիմվոլները թվային ուղղի ձախից կիսաբաց և աջից կի\-սա\-բաց ինտերվալների տոպոլոգիաներն են։

% 9.16
\item Ապացուցեք․ $(\R, \mapsto)$ և $(\R, \mapsfrom)$ տոպոլոգիական տարածությունները հոմեոմորֆ են։

\begin{hint}
Դիտարկեք $f: \R \to \R, \ f(x) = -x$ արտապատկերումը և օգտվեք \hyperref[ch9thm10]{թեորեմ 10}-ից։
\end{hint}

% 9.17
\item Ճի՞շտ է արդյոք, որ $(\R, \mapsto)$ և $(\R, \text{սովոր.})$ տարածությունները հոմեոմորֆ են։

\begin{hint}
Ապացուցեք, որ $[a, b)$ տեսքի ենթաբազմությունները միաժա\-մա\-նակ բաց և փակ ենթաբազմություններ են $(\R, \mapsto)$ տարածությունում, և օգտը- \linebreak վեք \hyperref[ch9thm10]{թեորեմ 10}-ից։
\end{hint}


% 9.18
\item Ունենք որևէ $f$, $g: X \to Y$ անընդհատ արտապատկերումներ, որտեղ $Y$-ը հաուսդորֆյան տարածություն է։ Դիտարկենք $F \subset X$ ենթաբազմությունը կազ\-մված այն բոլոր $x$ կետերից, որ $f(x) = g(x)$։ Ապացուցեք, որ $F$-ը փակ ենթա\-բազ\-մություն է $X$-ում։

\begin{hint}
Ցույց տվեք, որ $F$-ը պարունակում է իր բոլոր հպման կետերը։
\end{hint}

% 9.19
\item Դիցուք $f$, $g: X \to Y$ արտապատկերումները, $Y$ տարածությունը և ${F \subset X}$ ենթաբազմությունը այնպիսին են, ինչպես նախորդ խնդրում։ Ապացուցեք․ եթե $F$-ը նաև ամենուրեք խիտ է $X$-ում, ապա $f$ և $g$ արտապատկերումնե\-րը համընկնում են $X$-ի բոլոր կետերում։

\end{enumerate}


\end{document}